% RESEARCH_DRAFT_SHA256: 639d35c116ee8e766e6f442a87617cfd2a605e9ed21e5c264b81c15c482018ec
% Layout and source-reference wrappers expanded for the web renderer.
\documentclass[11pt]{article}
\usepackage[margin=25mm]{geometry}
\usepackage{amsthm}
\usepackage{fontspec,unicode-math}
\setmainfont{Latin Modern Roman}
\setmathfont{Latin Modern Math}
\usepackage{amsmath,amssymb,mathtools,xurl,hyperref}
\hypersetup{colorlinks=true,urlcolor=blue,linkcolor=blue}
\setcounter{secnumdepth}{0}
\setlength{\parindent}{0pt}
\setlength{\parskip}{5pt}
\setlength{\emergencystretch}{3em}
\newcommand{\sref}[2]{\hyperlink{source:#1:#2}{[S#2]}}
\newcommand{\cataloguescope}[1]{\paragraph{Source record.}#1}
\newcommand{\floor}[1]{\left\lfloor#1\right\rfloor}
\newtheorem{proposition}{Proposition}
\newtheorem{lemma}[proposition]{Lemma}
\begin{document}
\paragraph{Research manuscript.} OpenAI \texttt{gpt-6-astra}, reasoning \texttt{ultra}; 11 October 2026.
\paragraph{Preserved source body.} The complete catalogue mathematical text and sources follow with layout and source links expanded, including source qualifications and its historical preparation notes. New research begins at the section headed Research attempt.

% UnsolvedMath ID 30005630; source catalogue code OWR-14297738-003
\section{Training-Loss Implications for Shallow Neural Networks}\label{30005630}
{\small UnsolvedMath ID 30005630 · Computer Science}
\subsection{Definitions and mathematical statement}

\paragraph{Shared notation.}Standard number systems are $\mathbb N,\mathbb Z,\mathbb Q,\mathbb R,\mathbb C$. Specialized notation and index ranges follow the cited source and the definitions below.


A shallow scalar-output network conventionally has one hidden layer: $f(x)=b_0+\sum_{j=1}^m a_j\sigma(w_j^{\mathsf T}x+b_j)$ for a specified activation $\sigma$. Given training pairs $(x_i,y_i)$ and loss $\ell$, empirical training loss is $\widehat R(f)=n^{-1}\sum_{i=1}^n\ell(f(x_i),y_i)$. Population prediction risk under a distribution $P$ is $R_P(f)=\mathbb E_{(X,Y)\sim P}[\ell(f(X),Y)]$; generalization concerns the gap $R_P(f)-\widehat R(f)$ or a bound on $R_P(f)$. Small training loss alone does not specify a guarantee without assumptions on the sample, network class, and target distribution.
\paragraph{Record scope and qualifications.}
The catalogue title is retained as metadata. Named certificates, contracts, bounds or frameworks in that title are not asserted as results of the original source extract.
Source-record limitation: The supplied assessment flags an indeterminate training-loss prompt; data assumptions, loss, target prediction guarantee and network class are not specified.
\paragraph{Statement recorded in the supplied source.}
Determine what guarantees on prediction or generalization follow from small training loss for shallow neural networks. \href{https://doi.org/10.4171/owr/2023/41}{[S2]}
\paragraph{Source status.}
The supplied catalogue labels this source record open; that metadata was checked against the supplied snapshot on 2026-10-11. The label is preserved for provenance and is not an independent verification of present open status or mathematical correctness. The catalogue literature assessment is dated 2026-08-22.
\subsection{Short English statement}
The source asks about training-Loss Implications for Shallow Neural Networks.
\subsection{Sources}
\begin{itemize}
\item[S1] Ulam AI and the UnsolvedMath Contributors, UnsolvedMath: A Curated Collection of Open Mathematics Problems, record 30005630 (OWR-14297738-003). Catalogue curation and metadata: CC BY 4.0. The source extract is retained; definitions, English summaries and qualifications are editorial additions.
\url{https://huggingface.co/datasets/ulamai/UnsolvedMath}.
\item[S2] Mini-Workshop: Interpolation and Over-parameterization in Statistics and Machine Learning​ (2024), Oberwolfach Reports, DOI 10.4171/owr/2023/41, pp. 2359--2376
\url{https://doi.org/10.4171/owr/2023/41}.
\item Source-specific Oberwolfach Reports licences, including CC BY-SA 4.0 where applicable, are retained; see THIRD\_PARTY\_NOTICES.md.
\end{itemize}
\paragraph{Preparation and review.}
The original source and UnsolvedMath contributors retain attribution for the source extract. Codex (GPT-6; exact serving revision not exposed) assisted with definitions, English summaries, scope notes and formatting. No new human mathematical authorship or independently verified solution is asserted. Permanent source ID 30005630 and code OWR-14297738-003 checked against the complete supplied registry and current repository. Distinct upstream records, including repeated questions and broad or damaged extracts, are retained for complete coverage of the source CS category. A repeated record is not a new independent conjecture. Complete mathematical-equivalence and current-literature review has not been performed. Preparation and final review follow contribution prompts C and D, with source gaps disclosed. Structural and rendering checks do not establish a new theorem or open status.
\label{end:30005630}

\section{Research attempt}
\subsection{Literature, source context and scope}
The full original catalogue body and source qualifications are preserved above. The relevant primary workshop passages were read in the cached publisher report. Historical claims and subsequent attributed results are distinguished from the mathematical work below. A restricted example, standard lemma or known theorem is not counted as new progress on the entire original target. Missing assumptions and repeated or rhetorical extracts are not silently turned into different solved problems.

\subsection{A width-three network rules out a loss-only implication}
The primary passage discusses implications of different training losses and possible equivalence to exact interpolation. Its extraction does not specify a population-risk theorem. We test one natural but overstrong reading: zero training loss should by itself bound prediction risk of a shallow network.
Let $n\ge1$, let training inputs be independent uniform variables on $[0,1]$, and let every label be zero. Use squared loss and ReLU $\rho(z)=\max\{0,z\}$. Set $h=1/(2n)$ and, for any amplitude $A>0$, define the width-three network
\[
f(x)=\frac{2A}{h}\big[\rho(x)-2\rho(x-h/2)+\rho(x-h)\big].
\]
\begin{proposition}
With probability at least $1/2$ its empirical squared loss is zero, while its population risk is exactly $A^2h/3$.
\end{proposition}
\begin{proof}
The network vanishes outside $[0,h]$, rises linearly from zero to $A$ on $[0,h/2]$, and falls linearly to zero on $[h/2,h]$. The event that all inputs avoid $[0,h]$ has probability $(1-h)^n\ge1-nh=1/2$; the inequality follows from a union bound on any input entering the interval. On that event all predictions at training points vanish. The population squared loss is its squared integral, namely
\[
2\int_0^{h/2}(2Ax/h)^2\,dx=A^2h/3.
\]
\end{proof}
Choosing $A=\sqrt{3R/h}$ makes that risk any prescribed $R>0$ while keeping the same positive-probability zero-loss event. The construction uses only three hidden units but has unbounded weight scale. A bounded norm, Lipschitz constant, output range, or a particular algorithmic selection rule would change the question; none is supplied in the extracted implication. This is not a statement that a specified training algorithm must select this network.
\subsection{Verification and remaining gaps}
Piecewise network values and the exact squared integral were checked for rational amplitudes and widths, along with the no-hit probability. This provides a genuine obstruction to the loss-only inference for the cleaned question. It does not characterize the implicit bias induced by different losses, prove a neural-tangent equivalence, or resolve the source's broader research direction. Those require assumptions absent from the short record, which remain absent rather than being silently filled in.

\section{References and source attribution}
Primary workshop report: \url{https://doi.org/10.4171/owr/2023/41}. Relevant passages read: 2023/41 pp. 2369--2370, classification-versus-regression discussion and shallow-network training-loss question.

All retrievals were checked on 11 October 2026. Original question attribution remains in the preserved Sources section. External theorems are attributed and not claimed as new discoveries or fully verified merely by citation.


\paragraph{Integration review.} Codex (GPT-6; exact serving revision not exposed) reviewed the written arguments and expanded layout and source-reference wrappers. Repository metadata and the production web renderer were checked using Python 3.13 and JavaScript with MathJax 3.2.2. This model-assisted review is not independent human proof verification; no PDF compilation or formal proof checking is asserted.

\section{Completion Estimate}
\noindent\textbf{COMPLETION: 0\%}
No substantiated new progress toward the entire original target is claimed. The worked results and failed route are delimited above; a known result, illustrative model, or auxiliary obstruction is not counted as original-target completion.

\section{Final status and attribution}
\textbf{Attempt status: unresolved.} This describes this research attempt, not a new ruling on the source catalogue's status or on external literature claims.
The mathematics and this self-audit were generated by OpenAI \texttt{gpt-6-astra}, reasoning \texttt{ultra}, on 11 October 2026. Codex (GPT-6; exact serving revision not exposed) orchestrated the batch, read and reviewed its written mathematical arguments, and contributed formatting/repository review. No human mathematical authorship, independent human proof checking, or formal verification is claimed. The historical preparation paragraphs in the reproduced catalogue body concern that body only.

Contribution prompt A was applied to the complete definition and supplied source record, followed by the adversarial checks of prompt B. Local reproduction files are \texttt{batch025/checks.py}, \texttt{computations.json}, and the per-ID audit JSON. They were not public links at preparation time. Numerical checks are finite and do not prove unrestricted results. No LaTeX engine was found locally; compilation was not performed by this research agent.

\end{document}
